\documentclass[11pt]{article}
\usepackage[margin=1in]{geometry}
\usepackage{amsmath,amssymb,amsthm}
\usepackage[T1]{fontenc}
\usepackage{hyperref}
\hypersetup{colorlinks=true,linkcolor=blue,citecolor=blue,urlcolor=blue,
pdftitle={NP-completeness of the initial odd-ear prefix},
pdfauthor={Alec Kriebel},
pdfsubject={Variable-threshold NP-completeness for initial consecutive odd open ears},
pdfkeywords={open ear decomposition, factor-critical graph, NP-completeness, Set Cover, connected perfectly matchable induced subgraph}}
\newtheorem{theorem}{Theorem}
\newtheorem{lemma}[theorem]{Lemma}
\newtheorem{proposition}[theorem]{Proposition}
\newtheorem{corollary}[theorem]{Corollary}
\newcommand{\NP}{\mathsf{NP}}
\newcommand{\PP}{\mathsf{P}}
\newcommand{\cone}{\mathbin{*}}
\title{NP-completeness of the initial odd-ear prefix}
\author{Alec Kriebel\\
\small Independent researcher\\
\small ORCID: \href{https://orcid.org/0009-0001-9320-500X}{0009-0001-9320-500X}}
\date{22 October 2026\\\small Preprint, version 1.0.0}
\begin{document}
\maketitle
\begin{abstract}
For a finite 2-vertex-connected graph \(G\), let \(\beta(G)\) be the largest
number of consecutive odd ears at the beginning of an open ear decomposition,
counting its initial circuit and its one-edge ears. We prove that deciding
\(\beta(G)\ge t\), with \(t\) part of the input, is NP-complete even for simple
graphs with a universal vertex whose deletion leaves a connected bipartite
graph. This answers the variable-threshold complexity question posed by
Seb\H{o} and by Benchetrit--Seb\H{o}. The reduction is from explicitly encoded
Set Cover. It uses the classical factor-critical formulation of \(\beta\),
a cone identity, and a block-loss inequality that permits arbitrary partial
blocks. In particular, it does not assume an optimal subgraph is a union of
whole blocks.
\end{abstract}

\section{Question and conventions}
All graphs are finite and undirected. The general input graph is loopless,
may have parallel edges, and is 2-vertex-connected: it has at least three
vertices and remains connected after deleting any one vertex. A circuit is
a connected 2-regular subgraph. An \emph{open ear decomposition} consists of
an initial circuit followed by paths with distinct old endpoints and new
internal vertices. Each edge occurs in exactly one ear, and the ears
partition the edge set. Paths of one edge are allowed. An ear is odd if
its number of edges is odd. The initial circuit is also counted as an ear.

For such a graph, \(\beta(G)\) is the maximum length of the initial
consecutive odd prefix. Once the first even ear occurs, later odd ears do
not count. If no odd initial circuit is possible, \(\beta(G)=0\).
Seb\H{o}'s Question~4 in the 2014 Oberwolfach report
\cite[printed p.~2914]{Sebo} asks for the complexity of computing this
parameter. Benchetrit--Seb\H{o} give these precise open-ear conventions,
the factor-critical interpretation, and the polynomial-computability
question \cite[\S\S1.1,3.1, Question~3.3]{BS}.

\begin{theorem}\label{thm:main}
The language of pairs \((G,t)\) with \(G\) as above, \(t\) an integer in
binary, and \(\beta(G)\ge t\), is NP-complete. NP-hardness holds on simple
2-vertex-connected graphs with a specified universal vertex \(z\) such
that \(G-z\) is connected and bipartite.
\end{theorem}

Thus an exact polynomial-time algorithm for \(\beta\) would imply
\(\PP=\NP\). The threshold varies with the input. This does not assert
hardness at any fixed threshold, and is consistent with the fixed-parameter
result in \cite{BS}. It concerns neither closed ears nor the total number
of odd ears throughout a decomposition.

\section{The cone identity}
A graph \(H\) is \emph{factor-critical} if \(H-v\) has a perfect matching
for every \(v\in V(H)\). We use the classical theorem of Lov\'asz:
a 2-vertex-connected graph is factor-critical if and only if it has an
all-odd open ear decomposition \cite{Lovasz,BS}. This is an established
structural input, not a new result of this note.

\begin{lemma}\label{lem:fc}
For a 2-vertex-connected graph \(G\),
\[
 \beta(G)=
 \max_H \bigl(|E(H)|-|V(H)|+1\bigr),
 \tag{1}\label{eq:fc}
\]
where \(H\) ranges over 2-vertex-connected factor-critical subgraphs of
\(G\), with value zero if the family is empty. The maximum can be taken
over induced subgraphs.
\end{lemma}
\begin{proof}
The union of an all-odd prefix is 2-vertex-connected and factor-critical
by the classical theorem. Each ear contributes one to
\(|E|-|V|+1\), including the initial circuit. This proves one inequality.

Conversely, an eligible \(H\) has an all-odd open decomposition by the
same theorem. It extends to \(G\): if a component of \(G-V(H)\) had only
one distinct neighbor in \(H\), that neighbor would be a cut vertex of
\(G\). Hence any such component has two distinct old neighbors; a path
through the component between them supplies an open ear. Repeat until
all vertices are present, then add all unused edges as one-edge ears.
This also proves extension starting from any later intermediate union.
The resulting decomposition of \(G\) begins with all the ears of \(H\).

For a fixed vertex set, adding its missing edges preserves both
factor-criticality and 2-vertex-connectivity and can only increase the
quantity in (1). Thus induced \(H\) suffice.
\end{proof}

For a graph \(B\), its \emph{cone} \(z\cone B\) adds one new vertex \(z\)
and one edge from \(z\) to every vertex of \(B\).
\begin{lemma}\label{lem:cone}
If \(B\) is connected and bipartite with at least two vertices, then
\(G=z\cone B\) is 2-vertex-connected, and
\[
 \beta(G)=\max\{\,|E(B[S])|:
   S\ne\varnothing,\ B[S]\text{ is connected and has a perfect matching}\,\}.
 \tag{2}\label{eq:cone}
\]
\end{lemma}
\begin{proof}
Deleting \(z\) leaves connected \(B\); deleting any other vertex leaves
connectivity through \(z\). A 2-vertex-connected factor-critical graph
cannot be bipartite. Indeed its order is odd, its bipartition classes
are nonempty, and deleting a vertex from the smaller class leaves
unequal class sizes, which precludes a perfect matching. Thus every
eligible \(H\) in (1) contains \(z\). Make it induced and write
\(H=z\cone B[S]\).

Deleting \(z\) shows that 2-vertex-connectivity of \(H\) requires
\(B[S]\) connected, and factor-criticality requires a perfect matching
of \(B[S]\). Conversely, if \(B[S]\) is connected with a perfect matching
\(M\), then \(|S|\ge2\), and the cone is 2-vertex-connected. Deleting \(z\)
uses \(M\). After deleting \(u\in S\), remove its matching edge \(uv\)
and match \(v\) to \(z\) instead. This proves factor-criticality.
Finally,
\[
 |E(H)|-|V(H)|+1
 =\bigl(|E(B[S])|+|S|\bigr)-(|S|+1)+1=|E(B[S])|.
\]
Both directions give (2). Since \(B\) contains an edge, the feasible
family is nonempty.
\end{proof}

\section{The reduction}
We reduce from Set Cover, whose NP-completeness is classical
\cite{Karp}. An instance explicitly lists a universe
\(U=\{1,\ldots,m\}\), sets \(A_1,\ldots,A_n\), and a budget \(s\).
It asks whether at most \(s\) sets cover \(U\).

\paragraph{Preprocessing.}
Budgets below zero give a NO instance. With a nonnegative budget an
empty universe gives YES. For a nonempty universe, no sets, an uncovered
element, or budget zero give NO. If the union is \(U\) and \(s\ge n\),
the answer is YES. These tests are polynomial. Map any trivial YES to
a triangle with threshold~1 and any trivial NO to a triangle with
threshold~2. A triangle has \(\beta=1\) and is the cone over \(K_2\), so
\emph{every} output remains in the restricted class of
Theorem~\ref{thm:main}. It now suffices to handle
\[
 m,n\ge1,\qquad 1\le s\le n,\qquad \bigcup_{j=1}^n A_j=U.
\]
Empty or repeated sets are allowed in this nontrivial input.

Put
\[
 T=\sum_{j=1}^n(1+|A_j|),\qquad
 k=T+n+3,\qquad M=(m+1)k^2,\qquad F=M-ks.
 \tag{3}\label{eq:constants}
\]
Here \(k\) is a block size; the output threshold is \(F\).
We shall use
\[
 k>T,\qquad k>s,\qquad k^2>ks+T,\qquad F>0.
 \tag{4}\label{eq:gaps}
\]
Indeed \(k-s\ge T+3\), so \(k(k-s)>T\).

\paragraph{Construction.}
Take pairwise disjoint copies \(Q_0,\ldots,Q_m\) of \(K_{k,k}\), with
left and right classes \(L_i,R_i\). Choose a distinguished right vertex
\(q_i\in R_i\). For each \(A_j\), add a new left vertex \(c_j\), adjacent
to \(q_0\) and to each \(q_i\) with \(i\in A_j\). There are no other
interblock edges. Call this graph \(B\). It is simple and bipartite, with
\[
 |V(B)|=2k(m+1)+n,\qquad |E(B)|=M+T.
\]
Every connector \(c_j\) meets the root block \(Q_0\); every element block
has a connector neighbor since the sets cover \(U\). Thus \(B\) is
connected. Output
\[
       (G,t)=(z\cone B,F).                       \tag{5}
\]
It belongs to the claimed simple cone class.

\begin{proposition}\label{prop:equiv}
The Set Cover instance is YES if and only if \(\beta(G)\ge F\).
\end{proposition}
\begin{proof}
Suppose first that a cover of size at most \(s\) exists. Delete redundant
sets to obtain an inclusion-minimal cover indexed by \(J\), and put
\(r=|J|\le s\). Each \(j\in J\) has a \emph{private element} \(p(j)\):
an element of \(A_j\) belonging to no other selected set. Otherwise
\(A_j\) could be deleted. These private elements are distinct.

Select all block vertices and the \(r\) connectors indexed by \(J\),
but omit one left vertex from \(L_{p(j)}\) for each \(j\in J\).
Call the resulting vertex set \(S\). Match \(c_j\) to \(q_{p(j)}\).
In each such block, the remaining unmatched vertices form
\(K_{k-1,k-1}\); match them internally. Match all other blocks internally
as \(K_{k,k}\). Distinct private elements make these matches disjoint,
and they cover every selected vertex.

Each retained block is connected since \(k\ge2\). Every selected
connector meets the root block, and coverage links every element
block to one of those connectors. Thus \(B[S]\) is connected.
Exactly \(r\) left block vertices have been omitted, in distinct blocks,
losing exactly \(kr\) internal edges. All right vertices remain, so
every incidence of a chosen connector remains. Consequently
\[
 |E(B[S])|=M-kr+\sum_{j\in J}(1+|A_j|)
          \ge M-ks=F.
\]
Lemma~\ref{lem:cone} gives \(\beta(G)\ge F\).

Conversely, use that lemma to choose \(S\) for which \(B[S]\) is connected,
has a perfect matching, and has at least \(F\) edges. Let \(r\) be its
number of selected connectors. For every block, set
\[
 a_i=|L_i\setminus S|,\qquad b_i=|R_i\setminus S|,
 \qquad d_i=k(a_i+b_i)-a_i b_i.
\]
Here \(0\le a_i,b_i\le k\), and \(d_i\) is exactly the number of lost
internal edges of \(Q_i\). No whole-block assumption has been made.

Every block retains at least one internal edge. Otherwise at most
\(M-k^2\) internal edges remain and the total connector contribution
is at most \(T\), giving
\[
 |E(B[S])|\le M-k^2+T<M-ks=F
\]
by (4), a contradiction. In particular, both the root block and every
element block have selected internal edges.

The perfect matching forces the selected bipartition classes to have
equal size. All connectors are left vertices, so
\[
 \sum_i(k-a_i)+r=\sum_i(k-b_i),
 \qquad\text{hence}\qquad \sum_i(a_i-b_i)=r.
\]
The exact identity
\[
 d_i-k(a_i-b_i)=b_i(2k-a_i)\ge0                 \tag{6}\label{eq:loss}
\]
therefore gives \(\sum_i d_i\ge kr\). There are only \(T\) connector
edges in the whole graph. Thus
\[
 M-ks\le |E(B[S])|\le M-kr+T,\qquad k(r-s)\le T<k.
\]
Integrality implies \(r\le s\).

Finally, fix an element \(i\). Its block has a selected internal edge,
and so does the distinct root block. Since \(B[S]\) is connected, a
path from the element block to the root must leave \(Q_i\) through a
selected connector \(c_j\). Such an edge exists only when \(i\in A_j\).
The \(r\) selected sets therefore cover every element, within budget.
\end{proof}

\paragraph{Encoding size.}
Let \(L\) be the length of the explicit Set Cover input. After the union
test, \(m,n,T=O(L)\), and hence \(k=O(L)\). The construction lists
\(O(L^2)\) vertices and \(O(L^3)\) edges; the binary threshold \(F\)
has polynomial bit length. Renumbering the elements is polynomial.
No succinct multiplicities, weights, implicit graph, or parallel edges
are used. Proposition~\ref{prop:equiv} proves NP-hardness.

\paragraph{Membership in NP.}
For \(t>0\), a certificate supplies a subgraph \(H\subseteq G\) and
a perfect matching of \(H-v\) for each \(v\in V(H)\). Check that
\(|V(H)|\ge3\), that \(H\) and every \(H-v\) are connected, that all
matching edges belong to the required deletion graph and cover every
remaining vertex exactly once, and that
\(|E(H)|-|V(H)|+1\ge t\). These are polynomial checks with a polynomial
certificate. Lemma~\ref{lem:fc} proves the required equivalence.
For \(t\le0\) the answer is immediate on valid inputs. Invalid input
graphs may first be rejected in polynomial time. This completes the
proof of Theorem~\ref{thm:main}.

\section{Relation to earlier work and verification}
The cone identity optimizes the number of \emph{all induced edges} on
a \emph{connected} perfectly matchable vertex set. These two qualifications
matter. Sharifov's polynomial algorithm \cite{Sharifov} permits
disconnected perfectly matchable selections. For example, take three
\(K_{3,3}\) blocks and two left connectors, one joining the distinguished
root right vertex to each of the other blocks. All block vertices alone
give 27 induced edges and a perfect matching, but three components.
The connected optimum is 25: retaining all three blocks requires both
connectors, whose matching balance loses at least six block edges and
adds at most four connector edges. The bound is attained by omitting
one left vertex in each nonroot block and matching the connectors there.
If a block retains no internal edge, at most
\(27-9+4=22\) edges remain; with fewer than two connectors one cannot connect
all three retained blocks. Thus connectivity cannot simply be dropped.

Weighted Connected Matching \cite{Gomes} instead maximizes the weight
of the \emph{matching edges}; it does not maximize the number of all
edges induced by the endpoints. Its related hardness theorem is
therefore credited without identifying the two objectives.
Constrained ear decompositions \cite{HN} and joins in matroids
\cite{BBN} likewise concern different constraints or total-parity
parameters. The new ingredient here is the explicit Set Cover reduction
for (2), including the partial-block inequality (6).

The original question was stated as open in 2014 and 2015
\cite{Sebo,BS}. A bounded primary-source and publication-register
investigation through 22 October 2026 authenticated no earlier
publication of the full theorem above. This is not an exhaustive
priority or exclusive-firstness certification, and the question's
source and the cited dependencies are not represented as publications
of this reduction. Exact source versions and scope limits are recorded
in the accompanying \texttt{DEPENDENCIES.md} and \texttt{PRIORITY.md}.

The accompanying standard-library Python diagnostics test the original
complete-ear prefix objective on finite graphs against independent
matching and induced-subgraph enumerations, the exact omission
identity, explicit full-size reduction witnesses, preprocessing cases,
and falsifiable controls. These finite tests supplement the proof;
they do not establish its universal quantifiers or certify priority.
Reproduction instructions and exact output are included in the source
archive.

\paragraph{AI involvement and status.}
AI tools were extensively involved in mathematical research, proof
development and reconstruction, computational verification, literature
investigation, manuscript preparation, editing, and adversarial auditing.
Alec Kriebel is the named human author. This manuscript is a preprint
and has not undergone formal independent human peer review. Internal
AI reviews are not such peer review or proof-assistant verification.
Authored prose is licensed CC BY 4.0 and authored code is MIT licensed.

\begin{thebibliography}{9}
\bibitem{Sebo}
A.~Seb\H{o}, Problems about uniform covers, with tours and detours,
in \emph{Combinatorial Optimization}, Oberwolfach Reports 51/2014,
pp.~2912--2915, Question~4 on printed p.~2914.
\url{https://ems.press/content/serial-article-files/46541}.

\bibitem{BS}
Y.~Benchetrit and A.~Seb\H{o},
Ear-decompositions and the complexity of the matching polytope,
arXiv:1509.05586v1 (2015), \S\S1.1, 3.1--3.2.
\url{https://arxiv.org/abs/1509.05586v1}.

\bibitem{Lovasz}
L.~Lov\'asz, A note on factor-critical graphs,
\emph{Studia Sci. Math. Hungar.} \textbf{7} (1972), 279--280.
Classical theorem used in the precise form stated in \cite[Theorem~3.2]{BS}.

\bibitem{Karp}
R.~M.~Karp, Reducibility among combinatorial problems,
in R.~E.~Miller and J.~W.~Thatcher (eds.),
\emph{Complexity of Computer Computations}, Plenum, 1972, pp.~85--103.
\url{https://doi.org/10.1007/978-1-4684-2001-2_9}.

\bibitem{Sharifov}
F.~Sharifov, Perfectly matchable subgraph problem on a bipartite graph,
\emph{RAIRO Oper. Res.} \textbf{44} (2010), 27--42.
\url{https://doi.org/10.1051/ro/2010004}.

\bibitem{Gomes}
G.~C.~M.~Gomes, B.~P.~Masquio, P.~E.~D.~Pinto,
V.~F.~dos Santos, and J.~L.~Szwarcfiter,
Weighted connected matchings, \emph{Theoret. Comput. Sci.}
\textbf{1036} (2025), 115149.
\url{https://doi.org/10.1016/j.tcs.2025.115149}.
Inspected proof version: arXiv:2202.04746v1 (2022),
\url{https://arxiv.org/abs/2202.04746v1}; journal metadata authenticated,
journal body not accessed.

\bibitem{HN}
F.~Havet and N.~Nisse, Constrained ear decompositions in graphs and
digraphs, \emph{Discrete Math. Theor. Comput. Sci.}
\textbf{21}(4) (2019), article~3.
\url{https://doi.org/10.23638/DMTCS-21-4-3}.

\bibitem{BBN}
Y.~Bai, K.~B\'erczi, and C.~Nalam,
Joins and ear decompositions beyond graphic matroids,
arXiv:2608.01059v1 (2026).
\url{https://arxiv.org/abs/2608.01059v1}.
\end{thebibliography}
\end{document}
